\documentclass[11pt,a4paper]{article}

\usepackage[T1]{fontenc}
\usepackage[utf8]{inputenc}
\usepackage{lmodern}
\usepackage[margin=2.5cm]{geometry}
\usepackage{amsmath,amssymb,amsthm}
\usepackage{mathtools}
\usepackage{enumitem}
\usepackage{parskip}
\usepackage{setspace}
\usepackage[authoryear,round]{natbib}
\usepackage[hidelinks]{hyperref}
\usepackage[dvipsnames]{xcolor}

\setstretch{1.5}

% ---------------------------------------------------------------
% Theorem environments
% ---------------------------------------------------------------
\newtheorem{proposition}{Proposition}
\newtheorem{lemma}{Lemma}
\theoremstyle{definition}
\newtheorem{definition}{Definition}
\theoremstyle{remark}
\newtheorem{remark}{Remark}

% ---------------------------------------------------------------
% Notation
% ---------------------------------------------------------------
\newcommand{\F}{\mathcal{F}}
\newcommand{\G}{\mathcal{G}}
\newcommand{\C}{\mathcal{C}}
\newcommand{\E}{\mathbb{E}}
\renewcommand{\P}{\mathbb{P}}
\DeclareMathOperator{\logit}{logit}
\DeclareMathOperator{\Bin}{Bin}
\DeclareMathOperator{\Bern}{Bernoulli}

\title{Anytime-valid testing with Bayesian borrowing for platform trials}
 % add co-authors and affiliations
\date{\today}

\begin{document}
\maketitle

We propose a test for platform and multidimensional basket trials that uses Bayesian borrowing and controls the frequentist type I error at any sample size.
These trials evaluate targeted treatments in many patient subgroups under one master protocol.
We call each subgroup a basket and test it against a prespecified clinically relevant response rate.
Many baskets enrol few patients and some enrol many.
Borrowing information across baskets can increase power and inform effect sizes.
Bayesian hierarchical models borrow efficiently, but are not designed for frequentist hypothesis testing.
Decisions are based posterior probabilities and do not have formal type I error guarantees (see eg REFERENCES).

We construct an anytime valid test based on e-variables (REFERENCES).
The test allows interim analyses at any time and the choice to stop or continue after each analysis (optional stopping and continuation), while the type I error guarantee holds at every finite sample size.
The alternative in these likelihood ratios is a prior taken from a Bayesian borrowing model fitted to all baskets (see eg REFERENCES).
The prior can be updated before any interim analysis.
The type I error guarantee therefore holds whether or not the borrowing model is correct.

Link the previous paragraph (methodological benefits of using e-variables) to general concept of platform trials (require this kind of flexibilty) via a concrete example (DRUP study).

\section{Setup}\label{sec:setup}

We consider baskets arranged in a two-way grid with row levels $j = 1,\dots,J$ and column levels $l = 1,\dots,L$.
In our application, the rows are tumour types and the columns are biomarker groups \textcolor{orange}{True?}
A basket is $c = (j,l) \in \C \subseteq [J]\times[L]$.
We consider the sequence of patients $i = 1, 2, \dots$ ordered by the time at which their outcome is observed.
The $i$-th patient belongs to basket $c_i \in \C$ and has outcome $Y_i \in \{0,1\}$, where $Y_i = 1$ denotes a response.

Let $D_0$ denote external data (or clinical knowledge) fixed before the trial.
The trial history at time $t$ is the history of the whole platform:
\begin{equation}\label{eq:filtration}
  \F_t = \sigma\big( D_0,\; \{(c_i, Y_i) : i \le t\} \big), \qquad t = 0, 1, 2, \dots
\end{equation}
The basket of patient $t$ is known before the outcome $Y_t$ is observed.
We write
\begin{equation}\label{eq:pre-outcome}
  \G_t = \sigma(\F_{t-1}, c_t), \qquad \F_{t-1} \subseteq \G_t \subseteq \F_t,
\end{equation}
for the history just before $Y_t$ is observed.
We do not use a filtration of basket $c$ alone.
The borrowed prior depends on the other baskets and would not be measurable with respect to it.
Interim information on pending outcomes, such as early scans, is deliberately not in $\F$.

For each basket $c$, we state the null hypothesis conditionally on the platform history:
\begin{equation}\label{eq:hypotheses}
  H_{0,c}\colon\; \P\big(Y_i = 1 \mid \G_i\big) \le p_0 \quad \text{for all } i \text{ with } c_i = c .
\end{equation}
If patients are independent and basket $c$ has a constant response rate $p_c$, then $H_{0,c}$ reduces to $p_c \le p_0$.
The conditional form also allows drift in basket $c$.
It places no restriction on the other baskets.

Each basket answers its own clinical question: is the treatment active in this subgroup?
We dedicate a separate type I error budget $\alpha$ to every basket and apply no multiplicity correction across baskets.
For basket $c$ we require
\begin{equation}\label{eq:criterion}
  \sup_{P \in H_{0,c}} P\big( \text{reject } H_{0,c} \text{ at some interim analysis} \big) \le \alpha ,
\end{equation}
however many analyses are performed and whatever happens in the other baskets.

% ===============================================================
\section{Construction}\label{sec:construction}
% ===============================================================

An e-value for $H_0$ is a non-negative statistic $E$ with $\E_P[E] \le 1$ for all $P \in H_0$.
A non-negative supermartingale with initial value $1$ yields an e-value at every stopping time, and by Ville's inequality it exceeds $1/\alpha$ at some time with probability at most $\alpha$.

Basket $c$ is analysed at times $0 = \tau_{c,0} < \tau_{c,1} < \tau_{c,2} < \dots$
These analysis times are not specified in advance and may depend on the data observed so far.
Formally, they are stopping times with respect to $(\F_t)$.
Block $k$ of basket $c$ contains the outcomes of basket $c$ observed between two consecutive analyses.
We write $n_{c,k}$ for its number of outcomes and $y_{c,k}$ for its number of responses:
\begin{equation}\label{eq:blocks}
  B_{c,k} = \{ i : c_i = c,\; \tau_{c,k-1} < i \le \tau_{c,k} \}, \qquad
  n_{c,k} = \# B_{c,k}, \qquad
  y_{c,k} = \sum_{i \in B_{c,k}} Y_i .
\end{equation}

At analysis $k-1$ we choose a prior $\pi_{c,k-1}$ for the response rate of basket $c$.
It is a probability measure on $(p_0, 1]$ that is determined by the history $\F_{\tau_{c,k-1}}$.
It may use all platform data up to that time, including basket $c$'s own earlier blocks.
It is then frozen for block $k$.

The block factor is the likelihood ratio of the new block only, mixed over the frozen prior:
\begin{equation}\label{eq:factor}
  e_{c,k} = \int_{(p_0,1]} \Big(\frac{p}{p_0}\Big)^{y_{c,k}} \Big(\frac{1-p}{1-p_0}\Big)^{n_{c,k}-y_{c,k}} \pi_{c,k-1}(\mathrm dp).
\end{equation}
At each analysis, the evidence of the previous analysis is multiplied by the factor of the new block:
\begin{equation}\label{eq:update}
  E_c(\tau_{c,k}) = E_c(\tau_{c,k-1}) \cdot e_{c,k}, \qquad E_c(0) = 1 .
\end{equation}
We reject $H_{0,c}$ at the first analysis with $E_c(\tau_{c,k}) \ge 1/\alpha$.
The evidence depends on the data through the per-block counts and the priors frozen at each analysis.

\begin{proposition}\label{prop:validity}
Under $H_{0,c}$, for any choice of priors and analysis times that satisfies the conditions above,
\begin{equation}\label{eq:ville}
  P\big( \exists\, k : E_c(\tau_{c,k}) \ge 1/\alpha \big) \le \alpha .
\end{equation}
\end{proposition}

\begin{proof}
Fix $p>p_0$ and define the one-patient likelihood-ratio increment
\[
L_i(p)=\left(\frac{p}{p_0}\right)^{Y_i}
       \left(\frac{1-p}{1-p_0}\right)^{1-Y_i}
\]
when $c_i=c$, and set the increment equal to $1$ otherwise. Write
$q_i=\P(Y_i=1\mid\G_i)$. Under $H_{0,c}$, whenever $c_i=c$,
\[
\E[L_i(p)\mid\G_i]
=q_i\frac{p}{p_0}+(1-q_i)\frac{1-p}{1-p_0}\leq 1,
\]
because the expression is increasing in $q_i$ for $p>p_0$ and equals
$1$ at $q_i=p_0$. For patients in other baskets the increment is $1$.
Thus the product process for this fixed $p$ is a non-negative
supermartingale.

At time $\tau_{c,k-1}$, the prior $\pi_{c,k-1}$ is measurable with respect
to the observed history and is held fixed over the next block. Integrating
the fixed-$p$ supermartingales with respect to this probability measure
therefore gives a non-negative supermartingale for the block. Optional
sampling at the next analysis time and multiplication by the evidence
accumulated in previous blocks preserve the supermartingale property.
Ville's inequality gives the stated bound.
\end{proof}

\begin{remark}[Scope of the guarantee]
The guarantee is for each basket separately. It is not family-wise error
control for the event that at least one of many null baskets is rejected.
A platform-wide claim requires a separately specified multiplicity strategy.
Moreover, the conditional null in Equation~\ref{eq:hypotheses} is stronger
than a marginal-average response-rate null. The allocation of each patient
must be known before that patient's outcome is observed. The ordering and
handling of simultaneous outcome reports, and any outcome-dependent delay
in observation, must be specified so that the conditional null remains
meaningful for the sequence actually analysed.
\end{remark}

\begin{remark}[No maximum sample size]\label{rem:open}
Proposition~\ref{prop:validity} bounds the probability that $E_c$ reaches $1/\alpha$ at \emph{any} analysis of an unbounded sequence of patients.
A basket therefore needs no maximum sample size: it may continue until $H_{0,c}$ is rejected, or forever.
A planned maximum, a budget or a calendar end of the platform is optional stopping and does not affect the guarantee.
Without borrowing, or when borrowing from baskets whose data no longer change, the product of the block factors telescopes to the mixture likelihood ratio
$\int (p/p_0)^{y}\{(1-p)/(1-p_0)\}^{n-y}\,\pi_{c,0}(\mathrm dp)$ after $n$ patients with $y$ responses, where $\pi_{c,0}$ is the first restricted prior.
If $p_c>p_0$ is constant and $\pi_{c,0}$ has positive density near $p_c$, its logarithm grows like $n\,\mathrm{KL}(p_c\,\|\,p_0)$, so $H_{0,c}$ is rejected after finitely many patients with probability one: a test of power one.
Under $H_{0,c}$, however, the basket continues forever with probability at least $1-\alpha$.
An open-ended design therefore needs a futility rule.
\end{remark}

\begin{proposition}[Anytime-valid futility]\label{prop:futility}
Fix $p_1>p_0$ and consider the reversed hypothesis
\begin{equation}\label{eq:h1}
  H_{1,c}\colon\; \P\big(Y_i = 1 \mid \G_i\big) \ge p_1 \quad \text{for all } i \text{ with } c_i = c .
\end{equation}
Define $e^{f}_{c,k}$ and $E^{f}_c$ as in Equations~\eqref{eq:factor} and~\eqref{eq:update} with $p_0$ replaced by $p_1$, and with priors $\pi^{f}_{c,k-1}$ on $[0,p_1)$ that are determined by $\F_{\tau_{c,k-1}}$ and frozen for block $k$; for example, the borrowing posterior $\tilde\pi_{c,k-1}$ of Section~\ref{sec:model} restricted to $[0,p_1)$.
Stop basket $c$ for futility at the first analysis with $E^{f}_c(\tau_{c,k}) \ge 1/\beta$.
Under $H_{1,c}$, the probability that basket $c$ is ever stopped for futility is at most $\beta$.
\end{proposition}

\begin{proof}
As for Proposition~\ref{prop:validity}, with the inequality reversed.
For fixed $q<p_1$ and $c_i=c$,
$\E[(q/p_1)^{Y_i}\{(1-q)/(1-p_1)\}^{1-Y_i}\mid\G_i] = q_i\, q/p_1 + (1-q_i)(1-q)/(1-p_1)$
is decreasing in $q_i$ and equals $1$ at $q_i=p_1$, so it is at most $1$ under $H_{1,c}$.
\end{proof}

\begin{remark}[Open-ended design]\label{rem:open-design}
Stop basket $c$ at the first analysis with $E_c \ge 1/\alpha$ (reject $H_{0,c}$) or $E^{f}_c \ge 1/\beta$ (futility), without a maximum sample size.
Each rule is optional stopping for the other process, so $P(\text{reject } H_{0,c}) \le \alpha$ under $H_{0,c}$ and $P(\text{futility}) \le \beta$ under $H_{1,c}$ hold for the combined design, whether or not the borrowing model is correct.
Since $p_0<p_1$, every response rate lies above $p_0$ or below $p_1$.
If $p_c$ is constant and the priors concentrate around it, at least one of the two processes therefore grows, and the basket stops after finitely many patients with probability one.
In that case $P(\text{reject } H_{0,c}) \ge 1-\beta$ whenever $p_c \ge p_1$: a power guarantee without a maximum sample size.
Response rates between $p_0$ and $p_1$ may end in either decision.
\end{remark}

\begin{remark}[Computational approximation]
Equation~\ref{eq:restricted} defines the target restricted posterior.
The current implementation in \texttt{additional\_functions.R} represents it by normalized weights
$w_g$ on grid points $p_g>p_0$. The block factor is evaluated directly as the
finite sum of the likelihood ratios at those points. One-way posterior
weights use numerical quadrature; two-way weights use numerical grids and
Rao--Blackwellised MCMC. Thus the finite mixture is evaluated directly, but
the continuous posterior is approximated.

A finite Monte Carlo sample retained above $p_0$ is another admissible
empirical prior, provided at least one draw is retained and the resulting
weights are normalized. The earlier implementation in \texttt{Additional\_files/legacy/}
uses a moment-matched scaled Beta distribution and an analytic factor.
The new study retains a scaled-Beta benchmark represented on the same grid.
Each approximation defines its own alternative prior: support, normalization
and selection from past information preserve validity when the prior is
frozen before the next block. Numerical approximation can still change
power and borrowing. Any simulation randomness used to construct a prior
must also be independent of future trial outcomes.
\end{remark}

% ===============================================================
\section{Borrowing prior}\label{sec:model}
% ===============================================================

The proof uses only that the prior is frozen and lives on $(p_0,1]$.
Any borrowing model can supply it, and the model plays no role in error control.
One possible choice is a binomial logit EX/NEX model on the grid.
With $\theta_c = \logit(p_c)$ and cumulative counts $y_c \mid p_c \sim \Bin(n_c, p_c)$ at the time of the analysis, each basket is either exchangeable with an additive row and column structure (EX) or stands alone (NEX):
\begin{equation}\label{eq:exnex}
\begin{aligned}
  &\delta_c \sim \Bern(w_c), \qquad \theta_c \mid \delta_c = 0 \;\sim\; N(m_c, v_c), \\
  &\theta_c \mid \delta_c = 1 \;=\; \mu + a_j + b_l + \varepsilon_c, \quad \varepsilon_c \sim N(0, \sigma_\varepsilon^2), \\
  &\mu \sim N(m_\mu, s_\mu^2), \quad a_j \sim N(0, \sigma_a^2), \quad b_l \sim N(0, \sigma_b^2), \quad
  \sigma_a, \sigma_b, \sigma_\varepsilon \sim \text{half-normal}.
\end{aligned}
\end{equation}
Here $a_j$ is the row effect, $b_l$ the column effect and $\varepsilon_c$ the basket-specific interaction.
NEX protects against strong interactions, such as a treatment that is active in one tumour type but not in another.

Let $\tilde\pi_{c,k-1}$ be the posterior law of $p_c$ given the data in $\F_{\tau_{c,k-1}}$.
The prior for block $k$ is this posterior restricted to $(p_0, 1]$:
\begin{equation}\label{eq:restricted}
  \pi_{c,k-1}(A) = \frac{\tilde\pi_{c,k-1}\big(A \cap (p_0, 1]\big)}{\tilde\pi_{c,k-1}\big((p_0, 1]\big)}, \qquad A \subseteq [0,1] \text{ Borel}.
\end{equation}
One possible computation uses posterior draws
$p_c^{(1)}, \dots, p_c^{(H)}$. If $H'>0$ draws exceed $p_0$, the block
likelihood ratio can be averaged over those retained draws:
\begin{equation}\label{eq:mc}
  e_{c,k} \approx \frac{1}{H'} \sum_{h:\, p_c^{(h)} > p_0} \Big(\frac{p_c^{(h)}}{p_0}\Big)^{y_{c,k}} \Big(\frac{1-p_c^{(h)}}{1-p_0}\Big)^{n_{c,k} - y_{c,k}} .
\end{equation}

The implemented grid calculation instead uses all normalized restricted
grid weights, including the Rao--Blackwellised weights returned by the
two-way engine, as described in the computational remark above.



% \bibliographystyle{plainnat}
% \bibliography{refs}

\end{document}
